% SPDX-FileCopyrightText: 2026 Will Cook
% SPDX-License-Identifier: CC-BY-4.0
\documentclass[11pt]{article}
% SPDX-FileCopyrightText: 2026 Will Cook
% SPDX-License-Identifier: CC-BY-4.0
%
% Shared preamble for the Erdős Problem Notes: the short, standalone,
% problem-owned manuscripts covering the expansion library `ErdosProblems`.
%
% One note owns one Erdős problem.  Each note repeats whatever context its own
% reader needs, so that a reader who arrives at a single problem never has to
% assemble it from the gateway paper.  Deliberate repetition across notes is the
% design, not an oversight.
%
% Source links resolve against one pinned formal-source commit, declared once
% below.  `scripts/check_problem_note_sources.py` verifies that every authored
% (file, line, declaration) triple names that exact declaration in the pinned
% snapshot, so a link cannot silently rot as later work moves lines.

\usepackage{paper-house-style}
\usepackage{array}
\usepackage{booktabs}
\usepackage{enumitem}
\setlist{topsep=0.35em,itemsep=0.2em}
\usepackage[colorlinks=true,linkcolor=linkinternal,urlcolor=linkexternal,
  citecolor=linkinternal]{hyperref}
\hypersetup{bookmarksnumbered=true,bookmarksopen=true,bookmarksopenlevel=1,
  pdfauthor={Will Cook},pdflang={en-GB}}

% ---------------------------------------------------------------------------
% Pinned formal source.  \commit names the checkpoint every link resolves
% against; the checker reads that snapshot rather than the working tree, so the
% printed coordinates stay correct however later waves move declarations.
% ---------------------------------------------------------------------------
\newcommand{\commit}{e5fd26a6d1b1a346430205eab5385b4c9565d004}
\newcommand{\commitshort}{e5fd26a6d1b1}
\newcommand{\repobase}{https://github.com/wcook04/plectis-lean-erdos249-257/blob/\commit}
\newcommand{\sourceurl}{https://github.com/wcook04/plectis-lean-erdos249-257/tree/\commit}
\newcommand{\PX}{ErdosProblems}

% Declaration names stay inside link targets.  The printed label is either a
% spaced, lower-case reading of the declaration name (\lref) or a short authored
% phrase (\lword); neither prints a module path or a raw Lean identifier.
\ExplSyntaxOn
\cs_new_protected:Npn \noteleanlabel #1
  {
    \tl_set:Nx \l_tmpa_tl { \tl_to_str:n {#1} }
    \regex_replace_all:nnN { _+ } { \c{space} } \l_tmpa_tl
    \regex_replace_all:nnN
      { ([a-z0-9])([A-Z]) }
      { \1\c{space}\2 }
      \l_tmpa_tl
    \tl_set:Nx \l_tmpa_tl { \tl_lower_case:n { \tl_use:N \l_tmpa_tl } }
    \tl_use:N \l_tmpa_tl
  }
\ExplSyntaxOff
\newcommand{\leanlabel}[1]{\noteleanlabel{#1}}
\newcommand{\lref}[3]{\href{\repobase/\PX/#1\#L#2}{\leanlabel{#3}}}
\newcommand{\lword}[4]{\href{\repobase/\PX/#1\#L#2}{#4}}
\newcommand{\lloc}[2]{\href{\repobase/\PX/#1\#L#2}{Lean source}}

% Links into a sibling library, named repository-relative.  The expansion
% library is implicit in \lref/\lword, because that is where problem-owned
% work lands.  Several problems have their headline theorem in the older
% reviewed #249/#257 corpus, and a note that could not name that library
% would have to paraphrase a checked statement instead of linking it.
\newcommand{\mref}[3]{\href{\repobase/#1\#L#2}{\leanlabel{#3}}}
\newcommand{\mword}[4]{\href{\repobase/#1\#L#2}{#4}}
\newcommand{\mloc}[2]{\href{\repobase/#1\#L#2}{Lean source}}
\emergencystretch=2em

\newtheorem{theorem}{Theorem}[section]
\newtheorem{proposition}[theorem]{Proposition}
\newtheorem{lemma}[theorem]{Lemma}
\newtheorem{corollary}[theorem]{Corollary}
\theoremstyle{definition}
\newtheorem{definition}[theorem]{Definition}
\newtheorem{example}[theorem]{Example}
\newtheorem{problem}[theorem]{Problem}
\theoremstyle{remark}
\newtheorem*{remark}{Remark}

\newcommand{\N}{\mathbb{N}}
\newcommand{\Npos}{\mathbb{N}_{>0}}
\newcommand{\Z}{\mathbb{Z}}
\newcommand{\Q}{\mathbb{Q}}
\newcommand{\R}{\mathbb{R}}
\newcommand{\lcm}{\operatorname{lcm}}
\newcommand{\ph}{\varphi}

% The series line printed under every note's title block.
\newcommand{\seriesline}{Erd\H{o}s Problem Notes}

% Production provenance.  The wording is fixed by the corpus provenance
% contract and reproduced verbatim in every manuscript in this repository, so
% the papers cannot drift into different accounts of how they were made.
\newcommand{\noteprovenance}{\provenancenote{\emph{Authorship and AI use.} The
prose, formal proofs, and software in this version were drafted and revised by
large-language-model agents under Will Cook's direction.  Cook set the
objectives and acceptance criteria, reviewed and authorised the manuscript, and
is responsible for its contents.  The agents are production tools, not authors.
See \emph{Statements and declarations}.}}

% The status paragraph every note carries, in its own words, before any result.
% Kernel checking establishes the exact Lean proposition; it does not establish
% that the proposition is the interesting one, that it is new, or that the
% Erdős problem is closed.
\newcommand{\statusboundary}{%
  \emph{Status.}  The problem treated here is open, and this note does not
  close it.  Every statement below marked as checked is a proposition that the
  pinned Lean kernel accepts from the sources this note links to, with no
  \texttt{sorry}, no added axiom, and no unchecked evaluation.  That is a claim
  about the formal statement, not about its mathematical interest, its
  novelty, or the original problem.  The unresolved obligations are named
  exactly, in their own section, and none of the finite computations,
  reductions, or no-go results here removes one of them.}

\renewcommand{\commit}{99f4bf47422abbd8757cbb22b50ba079d764d3a7}
\renewcommand{\commitshort}{99f4bf47422a}
\hypersetup{
  pdftitle={Bases and integral relations for the k-kernel of Euler's totient},
  pdfsubject={Erdős Problem Notes: Problem 249},
  pdfkeywords={irrationality, Euler totient, k-regular sequences, 2-kernel, Lambert series, Möbius inversion, linear independence, Lean 4},
}

\newtheorem{equivform}[theorem]{Equivalent formulation}


\runninghead{Erd\H{o}s \#249: bases for totient sections}
\title{Bases and Integral Relations for the $k$-Kernel of Euler's Totient}
\subtitle{Exact finite-level ranks, with supplementary results on
Erd\H{o}s Problem~\#249}
\author{Will Cook}
\date{July 2026; revised 18 September 2026}

\begin{document}
\maketitle
\noteprovenance

\begin{abstract}
For integers $k\ge2$ and $e\ge1$, we give a basis for the rational span
of the totient sections $n\mapsto\varphi(k^jn+r)$ with
$0\le j\le e$ and $0\le r<k^j$. Its dimension is $k^e+1$.
Every section is an explicit integer multiple of a retained one, and
these reductions generate all integral relations, with unique
coefficients. The proof combines the local formula for Euler's totient
with a nonsingular evaluation matrix obtained from the Chinese remainder
theorem and Dirichlet's theorem. Supplementary results concern bounded
residue series and conditions for irrationality of
$S=\sum_{n\ge1}\varphi(n)2^{-n}$. They do not establish that $S$ is
irrational.
\end{abstract}


\section{A basis and all its relations}\label{sec:results}
We write $\ph$ for Euler's totient function and use
$\N=\{0,1,2,\ldots\}$, with $\ph(0)=0$. For $k\ge2$,
the $k$-kernel consists of the sequences
$n\mapsto\ph(k^jn+r)$ with $j\ge0$ and $0\le r<k^j$.
Some of these sections are scalar multiples of others: for example,
$\ph(4n+2)=\ph(2n+1)$ and $\ph(4n)=2\ph(2n)$. Thus the seven sections
through level two in base two are spanned by
\[
 n\longmapsto\ph(n),\quad\ph(2n),\quad\ph(2n+1),\quad
 \ph(4n+1),\quad\ph(4n+3).
\]
Are these five sequences independent, and do the analogous reductions
account for every relation in every base? The next theorem answers both
questions; Corollary~\ref{cor:integral-normal-form} gives the integral
relation module.

An integer sequence is $k$-regular if the $\Z$-module generated by its
$k$-kernel is finitely generated~\cite{allouche-shallit}. Coons proved that the totient sequence is
not $k$-regular for any $k\ge2$~\cite[Thm.~3.2]{coons}.
Martin's Theorem~1 already implies the affine independence used
below~\cite[Thm.~1]{martin-phi-inequalities}. The theorem below
identifies the finite-level basis and all integral relations.
\begin{theorem}[A basis through each finite level]
\label{thm:kkernelrank}
Let $k\ge2$ and $e\ge1$ be integers, write $F^{(k)}_{j,r}(n)=\ph(k^jn+r)$, and put
\[
 V_{k,e}=\operatorname{span}_{\Q}
 \{\,F^{(k)}_{j,r}:0\le j\le e,\ 0\le r<k^j\,\}.
\]
Then $\dim_{\Q}V_{k,e}=k^e+1$, and
\[
 \mathcal B_{k,e}
 =\{F^{(k)}_{0,0},F^{(k)}_{1,0}\}
 \cup\{\,F^{(k)}_{j,r}:1\le j\le e,\ 1\le r<k^j,\ k\nmid r\,\}
\]
is a basis.  Every omitted section reduces to a basis element by an explicit
scalar: $F^{(k)}_{j,0}=k^{j-1}F^{(k)}_{1,0}$ for $j\ge1$, and if $r=k^tu$ with
$t=\max\{s:k^s\mid r\}\ge1$ and $k\nmid u$, then
\[
 F^{(k)}_{j,r}=C_k(t,u)\,F^{(k)}_{j-t,u},
 \qquad
 C_k(t,u)=k^t\prod_{\substack{p\mid k\\ p\nmid u}}\Bigl(1-\tfrac1p\Bigr).
\]
\end{theorem}
\leannote{Lean: \lproof{ErdosProblems/Erdos249/PaperCompleteR8/FullKernelAssemblies.lean}{35}{displayed_all_base_kernel}.}
The condition on a retained positive residue is $k\nmid r$, not
$\gcd(k,r)=1$. Thus $r=2$ is retained when $k=6$. No primality or
squarefreeness assumption on the base is needed. The reduction scalar can
depend on the residue even for a fixed base: at $k=6$,
\[
 \ph(36n+6)=2\ph(6n+1),\qquad
 \ph(36n+12)=4\ph(6n+2).
\]
The second scalar differs because the argument $6n+2$ already contains
the prime $2$.

\subsection{Proof of the basis theorem}\label{sec:rank}
First consider spanning. The zero-residue reduction follows from
$\ph(k^jn)=k^{j-1}\ph(kn)$ for $j\ge1$: the prime divisors are unchanged
when $n>0$, and both sides vanish at $n=0$.
For a positive residue and a prime $p\mid k$ with $p\nmid u$,
the argument $k^{j-t}n+u$ is a $p$-adic unit. Multiplication by $k^t$
therefore contributes its power of $p$ and one factor $1-1/p$ to the
totient. When $p\mid u$, the argument already contains $p$, and only
the power contributes. This gives $C_k(t,u)$ and reduces every omitted
section to the displayed family.

For independence, we make an evaluation matrix diagonal and nonzero
modulo a prime. Each row requires just one primitive affine form to take
a prime value; auxiliary prime divisors make the other totients vanish
modulo the chosen prime. The two zero-residue sections are treated
separately at the end.

Consider $s$ forms $L_i(n)=a_in+b_i$, $1\le i\le s$,
with integer coefficients, positive slopes and
$a_ib_j\ne a_jb_i$ for $i\ne j$. We may restrict to a tail and shift
the variable so that every intercept is positive. Set $g_i=\gcd(a_i,b_i)$ and $A_i=L_i/g_i$;
the slope and intercept of $A_i$ are coprime. Choose an odd prime $\ell$ avoiding the slopes, all $g_i\ph(g_i)$, and the
nonzero cross determinants. In row $i$, choose distinct primes
$q_{ij}\equiv1\pmod\ell$ for $j\ne i$, outside the same finite exceptional
set, and impose
\[
 L_j(n)\equiv0\pmod{q_{ij}},\qquad A_i(n)\equiv2\pmod\ell.
\]
The Chinese remainder theorem combines these congruences. The progression for $A_i(n)$ has modulus
$(a_i/g_i)\ell\prod_{j\ne i}q_{ij}$ and is reduced. Primitivity gives
$\gcd(A_i(n),a_i/g_i)=1$ for every integer $n$, and the residue $2$
excludes $\ell$. If $q_{ij}$ divided $A_i(n)$, then it would divide both
$L_i(n)$ and $L_j(n)$, hence their nonzero cross determinant,
contrary to its choice.
Dirichlet's theorem supplies arbitrarily large primes $p=A_i(n)$ with
$p\nmid g_i$. Since $L_i(n)=g_ip$, the diagonal value is
$\ph(L_i(n))=\ph(g_i)(p-1)\not\equiv0\pmod\ell$.
Each off-diagonal value is divisible by $\ell$, since $q_{ij}-1$
divides the corresponding totient.
Choose such an evaluation point $n_i$ for each row $i$. The matrix
$\bigl(\ph(L_j(n_i))\bigr)_{i,j}$ is diagonal
modulo $\ell$, with every diagonal entry nonzero. Its determinant is
nonzero modulo $\ell$, hence is a nonzero integer. This proves rational
linear independence.

To handle the two proportional zero-residue forms, restrict to
$n=km+1$, where $\ph(kn)=\ph(k)\ph(n)$ by coprimality.
The positive-residue sections then give pairwise nonproportional forms
\[
 k^{j+1}m+(k^j+r),\qquad k\nmid r,
\]
and the remaining form is $km+1$. Proportionality between two of the
positive-residue forms would give $r/k^j=r'/k^{j'}$. If $j'>j$, then
$r'=k^{j'-j}r$ is divisible by $k$, contrary to the choice of $r'$;
if $j=j'$, then $r=r'$. None is proportional to $km+1$, since $r>0$.
The matrix therefore eliminates every positive-residue coefficient. If $A,B$ are the
two zero-residue coefficients, the restriction gives
$A+B\ph(k)=0$. Evaluation at $n=k$ gives $A+Bk=0$, since
$\ph(k^2)=k\ph(k)$. As $\ph(k)<k$, both coefficients vanish.
The reductions give spanning and the count is
\[
 2+\sum_{j=1}^e(k^j-k^{j-1})=k^e+1.
\]
This proves Theorem~\ref{thm:kkernelrank}. The linked formal statement gives
the same scalar reduction at level $j-t$, with its Euler product written out:
\href{https://github.com/wcook04/plectis-erdos/blob/f36a98bf3d3e6f65f1074e3b800e3293d5b8a51a/ErdosProblems/Erdos249/PaperCompleteR7/KernelIntegral.lean#L226}{reduction by the largest power of the base dividing the residue}.

The relation to the earlier CRT argument and the alternative proof from
Martin's theorem are discussed in Appendix~\ref{app:sources}.

\begin{corollary}[Dyadic basis]\label{res:basis}
The family
\[
 \{\ph_{0,0},\ph_{1,0}\}
 \cup\{\ph_{j,r}:j\ge1,\ 0<r<2^j,\ r\text{ odd}\},
 \qquad \ph_{j,r}(n)=\ph(2^jn+r),
\]
is a basis for the rational span of the full dyadic kernel. Every rational
relation is generated by the scalar reductions. The complete level-zero
truncation has dimension one.
\end{corollary}
\leannote{Lean: \lproof{ErdosProblems/Erdos249/PaperCompleteR8/FullKernelAssemblies.lean}{156}{displayed_full_dyadic_basis}.}
\begin{proof}
Every finite subset of the displayed family lies in $\mathcal B_{2,e}$
for some $e\ge1$, so the theorem gives independence. The reductions give
spanning.
At level zero the only section is $\varphi(n)$, which is nonzero since
$\varphi(1)=1$, so its span has dimension one.
\end{proof}

\begin{corollary}[Integral coordinates and all integral relations]\label{cor:integral-normal-form}
Let $k\ge2$ and $e\ge1$. The retained family is a $\Z$-basis of the
module generated by the sections through level $e$. Index the sections
by their level and residue, retaining distinct indices even when they
define equal sequences. For each omitted index $i$, write the scalar
reduction as $F_i=a_iF_{j(i)}$, where $j(i)$ is retained and $a_i$ is a
nonnegative integer. In the free abelian group with one generator $E_i$
for each of these indices, the vectors
\[
 R_i=E_i-a_iE_{j(i)}\qquad(i\text{ omitted})
\]
form a $\Z$-basis of the kernel of evaluation $E_i\mapsto F_i$.
In particular, every integral relation has a unique integral expression
in these elementary relations, and their rank is
$\sum_{j=1}^{e-1}k^j$.
\end{corollary}
\leannote{Lean: \lproof{ErdosProblems/Erdos249/PaperCompleteR8/KernelRelationBasis.lean}{398}{displayed_integral_normal_form}.}
\begin{proof}
The denominator $\prod_{p\mid k,\,p\nmid u}p$ in $C_k(t,u)$ divides $k$,
so the successive section reductions have integer scalars. They give
integral coordinates in the retained family; its rational independence
makes these coordinates unique over $\Z$ as well.

Now let $x=\sum_i x_iE_i$ evaluate to zero. For each omitted index $i$,
subtract $x_iR_i$. The result is supported on retained indices and still
evaluates to zero, so independence forces it to vanish. Thus the $R_i$
generate the integral relation module. Conversely, the coefficient at
an omitted index $i$ in $\sum_h c_hR_h$ is exactly $c_i$: each relation
has coefficient one at its own omitted coordinate and zero at all the
others. This proves independence and uniqueness without dividing by
any integer. There are
\[
 (1+k+\cdots+k^e)-(k^e+1)=\sum_{j=1}^{e-1}k^j
\]
omitted sections. The formal statement combines the integral coordinate
basis, these two-term relation vectors, and their exact rank.
This complete integral normal form is \href{https://github.com/wcook04/plectis-erdos/blob/25ef6245d15a47548c6926369ae8f1a0f0a14a80/ErdosProblems/Erdos249/PaperCompleteR8/KernelRelationBasis.lean#L398}{kernel-checked in Lean}.
\end{proof}

For example, at base $2$ and depth $2$, retain
$F_{0,0},F_{1,0},F_{1,1},F_{2,1},F_{2,3}$. The two omitted sections satisfy
$F_{2,0}=2F_{1,0}$ and $F_{2,2}=F_{1,1}$, so
\[
 E_{2,0}-2E_{1,0},\qquad E_{2,2}-E_{1,1}
\]
are an integral basis of all relations among the seven sections. The
rank calculation alone would give two; the basis also identifies every
relation and its unique coefficients.

The assertion concerns the integer combinations of the original sections,
not every integer-valued sequence in their rational span. For example,
at base $2$ and $e\ge2$, $\ph(4n+3)/2$ takes integer values but has
coefficient $1/2$ in the retained basis. The generated module is therefore
not saturated in the group of integer-valued sequences: an
integer-valued sequence outside the module can have a nonzero integer
multiple inside it.

\subsection{Relations with periodic coefficients}
Pairwise nonproportionality is essential here. For instance, $n+1$ and
$n+2$ satisfy the hypothesis. In contrast, $4n+2$ and $2n+1$ are
proportional, and $\ph(4n+2)=\ph(2n+1)$ gives a nontrivial relation.
Periodicity imposes no bound on the common period: it permits any finite
periodic pattern, because one can work on each residue class separately.

\begin{corollary}[Periodic coefficients]\label{cor:periodic-freezing}
Let $L_1,\ldots,L_s$ be pairwise nonproportional affine forms with integer
coefficients and positive slopes. If $w_1,\ldots,w_s$ are rational-valued
periodic sequences, then
\[
 \sum_i w_i(n)\ph(L_i(n))=0\quad\text{for all sufficiently large }n
 \quad\Longrightarrow\quad w_i(n)=0\quad\text{for every }i,n.
\]
\end{corollary}
\leannote{Lean: \lproof{ErdosProblems/Erdos249/PaperCompleteR20/PeriodicIntegerAffine.lean}{24}{periodic_freezing_integer_affine}.}
\begin{proof}
Write $L_i(n)=a_i n+b_i$ and choose a common positive period $Q$ for
the $w_i$. On a progression $n=r+Qt$, with $0\le r<Q$, their values
are the constants $w_i(r)$. The restricted affine forms have cross
determinants $Q(a_i b_j-a_j b_i)\ne0$. Shifting $t$ beyond the stated
threshold preserves these determinants and makes all arguments
positive. Affine independence then gives $w_i(r)=0$ for every $i$.
Doing this for each residue $r$ proves the conclusion at all indices,
by periodicity.
\end{proof}
Applied on each progression, Martin's Theorem~1 gives the same
conclusion~\cite[Thm.~1]{martin-phi-inequalities}.
The same argument preserves the basis after deletion of a finite prefix.
For the two sections with residue zero, large $n\equiv1\pmod k$ and large powers of $k$
replace the two evaluations used above.
\paragraph{Scope.}
This completes the basis argument. It describes relations among
coefficient subsequences, not the arithmetic nature of
$S=\sum_{n\ge1}\ph(n)2^{-n}$.

\bigskip\hrule\medskip
\section*{Supplementary results on the binary totient series}
The independent results below distinguish bounded-residue series,
finite exclusions for $S$, and conditions for irrationality. None is
needed for Theorem~\ref{thm:kkernelrank}; in particular, the conditions
in Sections~\ref{sec:carry-rank} and~\ref{sec:open} are not assumptions
of that theorem. Generating-function context and the formal-source guide
are in Appendix~\ref{app:sources}.

\section{Bounded residues and rationality}\label{sec:family}
After reducing $\ph(n)$ modulo a fixed integer, the coefficients are
bounded. A single nonzero coefficient between long zero blocks then
contradicts rationality by an elementary tail estimate.

\begin{theorem}[Residue series and dyadic observables]\label{res:residueseries}
For every $m\ge3$,
\[
 \sum_{n\ge1}\frac{\ph(n)\bmod m}{2^n}\notin\Q.
\]
For $k\ge1$ and $f:\Z/2^k\Z\to\Q$, the series
$\sum_{n\ge1}f(\ph(n)\bmod2^k)2^{-n}$ is rational exactly when $f$ is
constant on the even residue classes. If that constant is $c$, its value is
$3f(1)/4+c/4$.
\end{theorem}
\leannote{Lean: \lproof{ErdosProblems/Erdos249/PaperCompleteR7/RationalObservableClassification.lean}{277}{short_note_residue_theorem}.}

The restriction $m\ge3$ is necessary: modulo $2$ the only nonzero
coefficients occur at $n=1,2$, and the sum is $1/2+1/4=3/4$.
For the second assertion, only $f(1)$ and its values on the even classes
can affect the sum. The other odd classes are never attained by
$\varphi(n)$.

The next lemma applies, for example, to the indicator of the squares:
its coefficients are bounded, and sufficiently large squares have no other
square within any prescribed distance. An eventually periodic sequence
with infinitely many nonzero terms cannot satisfy the isolated-coefficient
hypothesis. Most importantly,
the bound $C$ must be fixed before the zero blocks are chosen; the
unbounded coefficients $\ph(n)$ do not satisfy this hypothesis.

\begin{lemma}[An isolated nonzero coefficient between long zero blocks]\label{lem:bounded-pulse}
Let $a_n\in\Z$ satisfy $|a_n|\le C$. Suppose that for arbitrarily large $L$
there is $N>L$ such that $a_N\ne0$ and
$a_{N+t}=0$ for $0<|t|\le L$. Then $\sum_{n\ge1}a_n2^{-n}$ is irrational.
\end{lemma}
\leannote{Lean: \lproof{ErdosProblems/Erdos249/PaperCompleteR7/PeriodicAndPulse.lean}{106}{bounded_isolated_pulse}.}
\begin{proof}
If the sum has denominator $q$, every scaled tail
$T_j=\sum_{t\ge1}a_{j+t}2^{-t}$ belongs to $q^{-1}\Z$.
Choose $L$ so large that $C2^{-L}<1/q$.
The right zero block gives $|T_N|\le C2^{-L}<1/q$, hence $T_N=0$.
The left zero block now gives
$T_{N-L-1}=a_N2^{-L-1}$, a nonzero element of $q^{-1}\Z$ with absolute
value less than $1/q$, a contradiction.
\end{proof}

\begin{proof}[Proof of Theorem~\ref{res:residueseries}]
Fix a unit $s\pmod m$ and a block length $L$. For each nonzero
$t\in[-L,L]$, choose a distinct prime $q_t\equiv1\pmod m$ larger than $L$
and outside the prime divisors of $m$. Prescribe
\[
 p\equiv s\pmod m,\qquad p\equiv-t\pmod{q_t}.
\]
This is a reduced progression, because each $q_t>|t|$. Dirichlet supplies
arbitrarily large primes $p$ in it. Then $m\mid\ph(p+t)$ for every
$0<|t|\le L$, whereas $\ph(p)\equiv s-1\pmod m$.

Now let $f:\Z/m\Z\to\Q$. Subtract $f(0)$ and clear the finitely many
denominators of its values. This produces a bounded integer sequence.
If $f(s-1)\ne f(0)$, the construction gives an isolated nonzero coefficient
between arbitrarily long zero blocks, so Lemma~\ref{lem:bounded-pulse}
proves irrationality. For the
least-residue map choose $s=-1$: its central residue is $m-2\ne0$ when
$m\ge3$. For $m=2^k$, every even residue $r$ has $r+1$ a unit, so a
nonconstant restriction to the even residues supplies such a nonzero coefficient.
Conversely, $\ph(n)$ is even for every $n\ge3$. Constancy on the even
classes makes all terms after $n=2$ constant, giving the stated rational
value without a recurrence hypothesis.
\end{proof}
There is an earlier coefficient-level obstruction: for each $m\ge3$,
$\varphi(n)\bmod m$ is not automatic in any integer base $k\ge2$
\cite[Cor.~4, p.~654]{yazdani2001}; see also
\cite[Thm.~3 and Ex.~4]{allouche-shallit-yassawi}.
Nonautomaticity does not prove irrationality of this sum: when $m>2$,
the residues are not binary digits. The preceding proof instead uses
the isolated coefficient and the bounded tail.
Erick Wong posted a proof that $\sum_{n\ge1}(\ph(n)\bmod k)k^{-n}$ is
irrational for every integer $k>2$, with the base equal to the modulus, as an
answer on Mathematics Stack Exchange on 29 March 2015~\cite{wong2015}.
Theorem~\ref{res:residueseries} fixes the base at $2$ for every modulus
$m\ge3$ and classifies the dyadic residue observables.

The preceding proof fixes its coefficient bound before choosing the zero blocks.
For $S$, the termwise condition $2^L\mid\ph(N+L)$, with $N+L>1$, forces
$2^L\le\ph(N+L)\le N+L-1$. Hence the tail envelope
$(N+L+2)2^{-L}>1$ cannot exclude any positive denominator.
The following criterion instead measures cancellation in a weighted prefix.

\section{Tail differences and finite residue tests}\label{sec:carry-rank}
Write
\[
 S=\sum_{n\ge1}\ph(n)2^{-n},\qquad
 R_N=\sum_{j\ge1}\ph(N+j)2^{-j},\qquad
 \Delta_h(N)=R_{N+h}-R_N.
\]
The identity
\[
 2^NS=\sum_{n=1}^N\ph(n)2^{N-n}+R_N
\]
has an integer finite sum. Subtracting the identities at $N+h$ and $N$
gives
\begin{equation}\label{eq:tail-phase}
 \Delta_h(N)\equiv2^N(2^h-1)S\pmod\Z.
\end{equation}
In particular, a nonintegral tail difference excludes every
rational value whose denominator divides $2^N(2^h-1)$.
\subsection{Certificates at arbitrarily large multiples of a shift}
For $h,N,L\in\N$, let
\[
 D_{h,N,L}=\sum_{j=0}^{L-1}
   \bigl(\ph(N+h+1+j)-\ph(N+1+j)\bigr)2^{L-1-j},
 \qquad B=N+h+L+2.
\]
All residues below are least nonnegative residues. Write $K(h,N,L)$ for
\[
 B<D_{h,N,L}\bmod2^L<2^L-B.
\]
The omitted terms form another difference of scaled tails:
\[
 2^L\Delta_h(N)-D_{h,N,L}=R_{N+h+L}-R_{N+L}.
\]
Since $0\le R_M\le\sum_{j\ge1}(M+j)2^{-j}=M+2$, the absolute
difference is bounded by the larger of the two tail bounds, not their sum.
Thus
\begin{equation}\label{eq:tail-error}
 |2^L\Delta_h(N)-D_{h,N,L}|\le B.
\end{equation}
A certificate $K(h,N,L)$ therefore forces $\Delta_h(N)\notin\Z$.
Conversely, write $\|x\|_{\R/\Z}=\min_{z\in\Z}|x-z|$ for the distance to
the nearest integer. If $\Delta_h(N)\notin\Z$, choose $L$ so that
$2^L\|\Delta_h(N)\|_{\R/\Z}>2B$. Such a choice exists: $h,N$ are fixed,
the distance is positive, and $B=N+h+L+2$ grows only linearly in $L$.
The error estimate then puts the residue of $D_{h,N,L}$ more than $B$
from either endpoint. Thus nonintegrality is equivalent to a certificate
at some unrestricted depth.

\begin{theorem}[Propagation of one nonintegral tail difference]\label{res:fulldepth}
Fix $d\ge1$ and $N\ge0$. If $\Delta_d(N)\notin\Z$, then every sufficiently
late pair $\{t,t+1\}$ contains an $m$ such that $K(md,N,md)$ holds.
Consequently
\[
 \exists t\ge1:\ K(td,N,td)
 \quad\Longleftrightarrow\quad\Delta_d(N)\notin\Z.
\]
Requiring this for every $d\ge1,N\ge0$ is equivalent to $S\notin\Q$.
\end{theorem}
\leannote{Lean: \lproof{ErdosProblems/Erdos249/PaperCompleteR7/ShortNoteAssemblies.lean}{35}{fullDepth_amplification}.}
For a fixed $d,N$, the hypothesis is just one nonintegral tail difference;
it does not assert irrationality of $S$. Requiring it for \emph{every}
$d,N$ is as strong as the original problem. The conclusion sets the
truncation depth equal to the shift, $L=md$, and loses at most one member
of each sufficiently late pair of successive multipliers.

\begin{proof}
Set $F_t=\Delta_{td}(N)$. Equation~\eqref{eq:tail-phase} yields
$F_{t+1}\equiv2^dF_t+F_1\pmod\Z$. If $K(td,N,td)$ fails,
\eqref{eq:tail-error} gives
\[
 \|F_t\|_{\R/\Z}\le\epsilon_t,
 \qquad \epsilon_t=2(N+2td+2)2^{-td}.
\]
Two adjacent failures would imply
$\|F_1\|_{\R/\Z}\le2^d\epsilon_t+\epsilon_{t+1}$, impossible for all sufficiently
large $t$ because the left side is positive and the right side tends to zero.
Conversely, if $F_1$ is integral, then every $F_t$ is integral by the same
recurrence, so no certificate exists. Finally, \eqref{eq:tail-phase} is
nonintegral for every $d,N$ when $S$ is irrational. If
$S=a/(2^cv)$ with $v$ odd, take $N=c$ and $d=\ph(v)$ to make it integral.
\end{proof}
For instance $D_{1,12,16}=-143140$, with residue $53468$ modulo $2^{16}$
and $B=31$, verifies the hypothesis at $d=1,N=12$. The theorem gives
infinitely many certificates with depth equal to the shift for this pair,
not for every $d,N$.

\noindent Formal sources: \lword{Erdos249/FullDepthRayAmplifier.lean}{364}{eventually_twoSyndetic_fullDepthKillMultipliers_of_seed}{a certificate in each sufficiently late pair of multipliers};
\lword{Erdos249/FullDepthRayAmplifier.lean}{406}{apFullDepthEscape_iff_irrational}{the equivalence after quantifying over all shifts and basepoints}.

\subsection{An equivalent residue condition for each hypothetical denominator}\label{sec:frontier}
For $c\ge0$, odd $v\ge1$ and a positive multiple $H$ of $\ph(v)$, put
\[
 M=\frac{2^H-1}{v},\qquad
 B_{H,c}=\sum_{j=0}^{H-1}\ph(c+1+j)2^{H-1-j}.
\]
Euler's theorem makes $M$ integral.

\begin{equivform}[A residue test modulo a divisor of $2^H-1$]\label{res:canonicalmersenne}
The series $S$ is irrational if and only if, for every $c\ge0$ and odd
$v\ge1$, some positive multiple $H$ of $\ph(v)$ satisfies
\begin{equation}\label{eq:canonical-gap}
 c+H+1<(-B_{H,c})\bmod M<M-(c+H+1).
\end{equation}
\end{equivform}
\begin{proof}
The identity
\[
 B_{H,c}=2^HR_c-R_{c+H}=M(vR_c)-\Delta_H(c)
\]
is exact. For $H\ge1$, $|\Delta_H(c)|<c+H+1$. Indeed $R_N\le N+1$
for $N\ge1$ and $0<R_0\le3/2$; the strict inequalities follow by subtracting
a positive tail. If $S=a/(2^cv)$, then $vR_c\in\Z$ and
$\Delta_H(c)\in\Z$. The residue of $-B_{H,c}$ is therefore within distance
less than $c+H+1$ of $0\pmod M$, so \eqref{eq:canonical-gap} fails.

If $S$ is irrational, fix $c,v$. Then $vR_c$ is nonintegral. As $H$ tends
to infinity through positive multiples of $\ph(v)$,
$\Delta_H(c)/M\to0$ and $(c+H+1)/M\to0$. Hence the fractional part of
$-B_{H,c}/M=-vR_c+\Delta_H(c)/M$ stays away from both endpoints of $[0,1]$,
which gives \eqref{eq:canonical-gap} for all sufficiently large such $H$.
\end{proof}
The condition is neither a routine size bound nor an additional assumption
known for a broad class containing $S$: with its universal quantifiers it
is equivalent to the unresolved irrationality assertion. A rational value
$a/(2^cv)$ fails the test for that same $c,v$, whereas irrationality gives
it for every sufficiently large allowed $H$. The proof does not find such
an $H$ without assuming irrationality.

The separate formal release states this equivalence using $vM=2^H-1$.
Its proof and the replay receipt are identified with the formal counterpart
to Lemma~\ref{lem:bounded-pulse}.
The finite residues also satisfy
\[
 \rho_{H,N+1,M}=
 \bigl(2\rho_{H,N,M}-\ph(N+H+1)+\ph(N+1)\bigr)\bmod M,
 \quad \rho_{H,N,M}=(-B_{H,N})\bmod M,
\]
which is an exact recurrence for searches directed at \eqref{eq:canonical-gap}.

\section{Information that does not force the gap}\label{sec:nogo}
\begin{enumerate}
\item\label{res:actualorbit}
\textbf{Tail differences at least common multiples.}  Let
$H_a=\lcm(1,\ldots,2^a)$. Then $S$ is irrational if and only if,
beyond every threshold $a_0$, some $a\ge a_0$ satisfies
$\Delta_{H_a}(H_a)\notin\Z$.  This is the
\mword{Erdos249257/TotientActualLcmOrbitNonintegrality.lean}{37}{irrational_totientSeries_iff_actualLcmOrbitNonintegralitySupply}{exact cofinal nonintegrality criterion}
for the actual LCM diagonal, not a quantitative
separation estimate and not a proof that such indices occur; the
\mword{Erdos249257/TotientActualLcmOrbitNonintegrality.lean}{53}{irrational_totientSeries_of_actualLcmOrbitNonintegralitySupply}{forward implication}
is recorded separately.

\item\label{res:actualsign}
\textbf{Positive tail differences need not be nonintegral.} For $a\ge8$ and
$J+(a+6)<2\cdot2^a$, one has $\Delta_{H_a}(H_a+J)>0$.
Suppose this difference is an integer. If a truncation depth $K$ satisfies
\[
 J+K+(a+6)<2\cdot2^a,
 \qquad 2H_a+J+K+2<2^K,
\]
then
\[
 D_{H_a,H_a+J,K}\bmod2^K
   =2^K-\Delta_{H_a}(H_a+J+K)
\]
lies strictly between $2^K-(2H_a+J+K+2)$ and $2^K$.
The later tail difference is a positive integer; its \emph{negative} is
the representative near zero. Thus positivity determines which endpoint
contains the residue, but does not exclude an integral tail difference.
Both displayed bounds on $K$ are required for this conclusion.
{\footnotesize\emph{Checked:}
\mword{Erdos249257/TotientActualLcmOrbitSign.lean}{39}{actualLcmTailDiff_shift_pos}{positive tail differences},
\mword{Erdos249257/TotientActualLcmOrbitSign.lean}{172}{actualLcm_trueEndpointSurvivor_neg}{negative endpoint representative},
\mword{Erdos249257/TotientActualLcmOrbitSign.lean}{211}{actualLcm_integral_forces_topEdgeResidue}{residue near the upper endpoint}.}

\item\label{res:factorideal}
\textbf{Divisibility and finite linear combinations do not suffice.}
A dyadic coboundary is a sequence of the form $2u_j-u_{j+1}$; its weighted
sum telescopes. For every $t\ge3$, with $H=\lcm(1,\ldots,t)$, the linked
source constructs a nonzero comparison sequence of this form. Its
coefficients reproduce the actual totient differences at $t-2$
prescribed indices and are divisible by $\varphi(H)$. Every finite integer
linear combination of forward shifts preserves the telescoping identity and
divisibility by $\varphi(H)$ and the corresponding lower-level factors.
The bounds remain uniform in the index, but are multiplied by the sum
of the absolute values of the combination's coefficients.
A variant preserves sparse prescribed values along entire arithmetic
progressions. These properties alone do not force a contradiction: the
comparison need not consist of actual totient differences. The theorem
concerns linear combinations of shifts, not nonlinear operations, and
does not supply certificates at arbitrarily large scales.
{\footnotesize\emph{Checked:}
\mword{Erdos249257/LcmFactorIdealPulseObstruction.lean}{798}{lcm_factorIdeal_finiteRank_shiftAlgebra_not_sufficient}{obstruction for finite integer combinations of shifts}.}
\end{enumerate}

If $S=a/v$ with integers $a$ and $v\ge1$, the scaled tails $u_N=vR_N$
are integers and satisfy
\[
 u_{N+1}=2u_N-v\ph(N+1),\qquad 0\le u_N\le v(N+2).
\]
In particular $u_N/2^N\to0$, the condition called temperedness in the
formal sources. Its dyadic sections $n\mapsto u_{2^jn+r}$ have rational rank at
least $2^e-1$ through level $e$ and a common eventual period modulo $v$.
These properties concern the same carry and are compatible; neither
supplies a contradiction.

\noindent Formal sources: \mword{Erdos249257/TotientTailCarryPeriod.lean}{224}{not_irrational_totientSeries_implies_mod_period_and_unbounded_rank}{same-carry rank and modular-period theorem}.

The linked formal source constructs a rational comparison sequence with
$c(n)=\ph(n)$ for odd $n$, $|c(n)-\ph(n)|\le2$ for even $n$, $0\le c(n)\le n$,
and
\[
 \sum_{n\ge1}c(n)2^{-n}=\frac54.
\]
This comparison sequence has an integral tempered carry with dyadic rank
at least $2^e-1$ at every level. Thus the same lower bound for a hypothetical totient carry is
compatible with rationality. For one fixed prime $p$ and $g:\N\to\Z$, the exact laws
$g(pn)=p g(n)$ when $p\mid n$ and $g(pn)=(p-1)g(n)$ otherwise, together
with $g-\ph=o(n)$, force $g=\ph$: a nonzero defect keeps a nonzero ratio
to the argument along $p^j n$. The comparison sequence fails these dilation laws. The unchanged odd-residue sections
also give coefficient rank at least $2^e-1$ at every level. This is a
lower bound, not an exact formula for the full coefficient rank.

A different obstruction concerns rational approximation. Let $\mu$ be the
M\"obius function and put
\[
 \Theta_2=\sum_{d\ge1}\frac{\mu(d)}{(2^d-1)^2}=S-\frac12,
 \qquad
 Q(e,Y)=\frac{\bigl(\sum_{d=1}^{Y}\mu(d)/(2^d-1)^{e+2}\bigr)^2}
 {\sum_{d=1}^{Y}\mu(d)/(2^d-1)^{2e+2}}.
\]
The divisor-convolution identity $\ph=\mu*\mathrm{Id}$ gives
$S=\sum_{d\ge1}\mu(d)2^d/(2^d-1)^2$.
Now $2^d=(2^d-1)+1$ and
$\sum_{d\ge1}\mu(d)/(2^d-1)=1/2$, the latter following from
$\sum_{d\mid n}\mu(d)=0$ for $n>1$ and $1$ for $n=1$.
These identities give the displayed value of $\Theta_2$; all the sums are
absolutely convergent. Amiram Eldar posted the corresponding formula
$S=\tfrac12+\sum_{d\ge1}\mu(d)/(2^d-1)^2$, together with its coprimality
interpretation, to OEIS A256936 on 15 March 2026~\cite{eldar2026oeis}, and
Steve Fan posted the formula on the problem's forum thread on
16 May 2026~\cite{fan2026totient}. The following exact extremum is
stated as the rank-one lower-bound theorem in the long record; its proof is
\lword{Erdos249/RankOneSharpFloor.lean}{534}{rankOneSubrankQuotient_eq_one_five_iff}{the sharp lower-bound argument}
in the linked source.

\begin{theorem}[A lower bound for the rank-one quotients]\label{res:rankonefloor}
For $e\ge1$ and $Y\ge4$, the denominator of $Q(e,Y)$ is positive, and the
unique minimiser is $(e,Y)=(1,5)$. Every admissible quotient and every
nonempty finite positive weighted average of such quotients satisfies
\[
 Q-\Theta_2>\frac{21}{320}.
\]
The uniform bound $Q-\Theta_2>1/15$ is false already at $(1,5)$.
\end{theorem}
\leannote{Lean: \lproof{ErdosProblems/Erdos249/RankOneSharpFloor.lean}{25}{rankOne_denominator_pos}, \lproof{ErdosProblems/Erdos249/RankOneSharpFloor.lean}{530}{rankOneSubrankQuotient_ge_one_five}, \lproof{ErdosProblems/Erdos249/RankOneSharpFloor.lean}{546}{rankOneSubrankQuotient_eq_one_five_iff}, \lproof{ErdosProblems/Erdos249/RankOneSharpFloor.lean}{636}{rankOneSubrankQuotient_sub_theta_two_gt_twentyOne_div_threeTwenty}, and 2 further declarations in the \papersectionlink{erdos249-totient-reasoning-surface.pdf}{coverage}{coverage section of the companion record}.}
The unique minimiser and the
\lword{Erdos249/RankOneSharpFloor.lean}{624}{rankOneSubrankQuotient_sub_theta_two_gt_twentyOne_div_threeTwenty}{uniform lower bound $21/320$}
are in the linked source; the $1/15$ bound
\lword{Erdos249/RankOneSharpFloor.lean}{645}{not_forall_rankOneSubrankQuotient_sub_theta_two_gt_one_div_fifteen}{already fails at $(1,5)$}.
This theorem treats averages after the quotients are formed. Signed
coefficients, coupling before the quotient, and higher-rank kernels require
separate arguments. For rational approximants $p_j/q_j$ in lowest terms, with $q_j>0$, a sufficient arithmetic target is
\[
 0<q_j\left|\Theta_2-p_j/q_j\right|\longrightarrow0.
\]
Convergence alone does not give this estimate. The elementary dyadic
prefix bounds give only $q_j|S-p_j/q_j|\le N+2$ at depth $N$, not a
quantity tending to zero. A prescribed approximation family therefore
needs a sharper error or reduced-denominator bound. When approximants
may be chosen freely, continued-fraction convergents show that the
criterion is equivalent to irrationality, just as
\eqref{eq:canonical-gap} is. The theorem above excludes the displayed
rank-one quotients and their positive averages, not other approximation
families.

\section{A remaining exponential-sum estimate}\label{sec:open}
A sufficient analytic target is a bound on an explicit finite sum:
for every $h\ge1$ and every $X_0$, find integers $X\ge\max(X_0,1)$ and
$L\ge h$ such that
\[
 16(2X+h+L+2)\le2^L,\qquad
 \sum_{N=X}^{2X-1}\cos\!\left(\frac{2\pi D_{h,N,L}}{2^L}\right)
                                      \le\frac9{10}X.
\]
Here $D_{h,N,L}$ is the integer window sum from Section~\ref{sec:carry-rank}.
The room inequality makes the truncation error small, whereas the
cosine bound forces at least one residue away from the integer lattice.
The long record proves that these quantified inequalities would imply
irrationality. The modulus grows with $L$, and the coefficients inside
the phase are totient differences; ordinary information about the
distribution of prime factors does not give this estimate.

One proposed approach separates a single totient value from each phase.
The following definitions specify the four estimates it would require.
Fix integers $h,s\ge1$, $L\ge s+h$, $X\ge1$ and $0<\eta<1$. Put
$t=L-s+1$ and $E_N=\exp(2\pi iD_{h,N,L}/2^L)$ for $X\le N<2X$.
Let $p_N$ be the largest prime factor of $N+t$ and $m_N=(N+t)/p_N$.
Define
\[
 \begin{aligned}
 \mathcal A&=\{N\in[X,2X):
       m_N\le\lfloor\sqrt X/2\rfloor,\quad p_N>2\lfloor\sqrt X\rfloor\},\\
 \mathcal G&=\{N\in\mathcal A:\ph(m_N)\ge\eta m_N\}.
 \end{aligned}
\]
All index sets here contain integers. On $\mathcal A$, the inequalities
ensure $p_N>m_N\ge1$, so $\ph(N+t)=\ph(m_N)(p_N-1)$.
The coefficient of this totient value in $D_{h,N,L}/2^L$ is
$(2^h-1)/2^t$. Set
\[
 z_N=\exp\!\left(\frac{2\pi i(2^h-1)\ph(m_N)(p_N-1)}{2^t}\right),
 \qquad w_N=E_N/z_N \quad(N\in\mathcal A),
\]
and let $\bar z_m$ be the average of $z_N$ over
$\{N\in\mathcal A:m_N=m\}$, with value zero for an empty set.
The proposed bounds are
\[
 \begin{aligned}
 \operatorname{Re}\sum_{N\in\mathcal G}w_N(z_N-\bar z_{m_N})
       &\le\tfrac{14}{25}X,\\
 \left\lVert\sum_{N\in\mathcal G}w_N\bar z_{m_N}\right\rVert
       &\le\tfrac1{100}X,\\
 \left\lVert\sum_{N\in\mathcal A\smallsetminus\mathcal G}E_N\right\rVert
       &\le\tfrac1{100}X,\\
 \left\lVert\sum_{\substack{X\le N<2X\\N\notin\mathcal A}}E_N\right\rVert
       &\le\tfrac8{25}X.
 \end{aligned}
\]
Before taking real parts or norms, these four sums add exactly to
$\sum_{N=X}^{2X-1}E_N$: partition the indices and use $w_Nz_N=E_N$.
Subtracting the group mean is an identity, not an independence assumption.
The four constants add to $9/10$, giving the preceding real-part bound;
they do not give a $21/25$ bound on the norm of the whole sum. The allocations
are sufficient margins, not claimed optimal constants. In particular,
$1-\log2<8/25$ leaves room for the smooth-number estimate used below.
\begin{problem}[Four estimates for the first-harmonic sum]
\label{prob:firstharmonic}
For each $h\ge1$, do there exist $s\ge1$ and $0<\eta<1$ such that, for every
$X_0$, there are $X,L$ with
\[
 \max(X_0,1)\le X,\qquad h\le L-s,\qquad
 16(2X+h+L+2)\le2^L,
\]
for which those four bounds hold?
\end{problem}
The long record gives the same decomposition, and the
\mword{Erdos249257/FirstHarmonicPivot.lean}{543}{PivotBudgetAt}{four simultaneous inequalities}
are defined in the linked source. All four bounds must hold on the \emph{same} block
with the same parameters. The permitted blocks need only be arbitrarily
large, but $L$ must also make the truncation error small enough for the
residue test. The long record proves the unassigned and excluded-cofactor
bounds at minimal admissible depth, using smooth-number asymptotics and
the distribution of $\ph(m)/m$. An elementary mean-value estimate also
gives the explicit sufficient cutoff $\eta=1/1000$ for the
excluded-cofactor bound. The two remaining estimates concern the group
means and their weighted, mean-subtracted phases. They must hold at the
same depth and for the same choice of $\eta$; the existence of a cutoff
that controls the excluded indices does not establish either of them.

Balasubramanian, Giri and Srivastav estimate the finite sums
$\sum_{1\le n\le x}F(n)G(n-h)$ for $F=f*1$, $G=g*1$ with decaying
$f,g$, setting $F(n)=G(n)=0$ for $n\le0$. Taking
$f(n)=g(n)=\mu(n)/n$ gives $F(n)=G(n)=\ph(n)/n$ for $n\ge1$;
the estimate is uniform for $|h|\le x/2$~\cite[Thm.~2.2]{bgs2016}.
Their partial-summation extension allows differentiable weights with
bounded derivative on the summation interval~\cite[Rem.~2.9]{bgs2016}.
The missing input here is a useful error bound for the present phase
and residual weight, not shift uniformity or partial summation itself.
\medskip
\noindent\textbf{Acknowledgement.} I thank Wouter van Doorn for advice on
writing for a first-time reader and on explaining the strength of a
hypothesis. His advice concerned a different note; it was not a mathematical
review or endorsement of the results presented here.

\appendix
\section{Further context and formal sources}\label{app:sources}

\subsection*{Affine independence in the earlier literature}
The separation method has an earlier source: Yazdani's proof of
Theorem~2, explicitly credited there to Shallit, uses the Chinese remainder
theorem and Dirichlet's theorem to distinguish kernel sections modulo a
fixed modulus~\cite[Thm.~2 and proof, pp.~652--653]{yazdani2001}.
For a later account see Allouche, Shallit and Yassawi~\cite[Thm.~3 and
Rem.~7]{allouche-shallit-yassawi}. Here one separates a chosen column
from every other column of a finite evaluation matrix. Nonsingularity,
not merely distinctness of infinitely many sections, gives linear
independence.

Martin's Theorem~1 gives another proof of the affine independence
conclusion~\cite[Thm.~1]{martin-phi-inequalities}. For any ordering of the nonproportional forms and any
$C>0$, their successive totient ratios exceed $C$ on a set of positive
lower density. In a proposed relation $\sum_i c_i\ph(L_i(n))=0$,
put a nonzero coefficient first and choose
$C>\max(1,\sum_{i>1}|c_i|/|c_1|)$. At the resulting values of $n$,
the first term has larger absolute value than the sum of all the
others, a contradiction. The proof in Section~\ref{sec:rank} instead constructs a finite
nonsingularity certificate for each chosen family.

\subsection*{Regularity, Mahler equations and section maps}
The totient is not $k$-regular, in the sense of Allouche and
Shallit~\cite{allouche-shallit}, for any $k\ge2$: Coons proved that the
$\Z$-module generated by its $k$-kernel is not finitely generated, by counting
the poles of $\zeta(s-1)/\zeta(s)$ in the strip
$0\le\operatorname{Re}s\le1$~\cite[Cor.~3.1, Thm.~3.2]{coons}. For
integer-valued sequences this is equivalent to infinite dimension of the
rational span: if the span has dimension $d$, some $d$ evaluations map it
injectively to $\Q^d$ and the module into $\Z^d$ (compare the evaluation
argument in~\cite[App.~A]{coons-jsr}). Theorem~\ref{thm:kkernelrank}
therefore recovers Coons's nonregularity theorem~\cite[Thm.~3.2]{coons}
and gives the dimension at every finite level.

The generating-function conclusion concerns functional equations, not
just the dimensions of section spans. A $k$-Mahler equation for a formal
power series $F$ has the form
\[
 \sum_{i=0}^{d}P_i(z)F(z^{k^i})=0,
 \qquad d\ge0,\quad P_i(z)\in\Q[z],\quad P_0\ne0.
\]
The height criterion of Adamczewski, Bell and Smertnig already implies
that an integer-coefficient $k$-Mahler series with polynomially bounded
coefficients is $k$-regular~\cite[Thm.~1.2(b), p.~2528]{adamczewski-bell-smertnig2023}.
Indeed, the logarithmic height of an integer $a$ is
$\log\max(1,|a|)$. Since $0\le\ph(n)\le n$, nonregularity rules out
such an equation for $F_\varphi(z)=\sum_{n\ge1}\varphi(n)z^n$ in every
base $k\ge2$; no multiplicativity hypothesis is needed for this deduction.
The multiplicative classification in Bell and Smertnig's later preprint
also gives this conclusion~\cite[Thm.~1.3]{bell-smertnig2026}.
Bell, Bruin and Coons give related classifications for algebraic or
$D$-finite generating functions~\cite[Thms.~1.5--1.6]{bell-bruin-coons};
$D$-finite means satisfying a linear differential equation with
polynomial coefficients.
These theorems constrain the whole generating function. They do not
by themselves determine the arithmetic nature of its single value
$F_\varphi(1/2)$.

For example, the next dyadic level has rank $9$ rather than $15$:
$\ph(8n+2)=\ph(4n+1)$, $\ph(8n+4)=2\ph(2n+1)$ and
$\ph(8n+6)=\ph(4n+3)$ are its positive-residue reductions.

For a $k$-regular sequence that is not eventually zero, its growth
exponent is the base-$k$ logarithm of the joint spectral radius of the
section matrices in a basis of the kernel span. A redundant spanning
family gives only an upper bound~\cite[Thm.~1, Prop.~4, Cor.~7]{coons-jsr}.
This finite-dimensional theorem does not apply to $\ph$. Its section
maps nevertheless have the explicit form below. For $0\le i<k$, put $T_i g(n)=g(kn+i)$. These are the
coefficient versions of the Cartier operators, and
\[
 T_iF_{j,r}=F_{j+1,k^ji+r},\qquad
 V_{k,e+1}=V_{k,e}+\sum_{i=0}^{k-1}T_iV_{k,e}.
\]
Write each residue at level $j+1$ as $k^ji+r$, $0\le r<k^j$, to
obtain the second equality. The strict dimension increase therefore
prevents all $T_i$ from preserving $V_{k,e}$. A retained positive-residue
section maps to another such section, since $k\nmid(k^ji+r)$ when
$j\ge1$ and $k\nmid r$.
The two zero-residue generators need separate treatment:
$T_iF_{0,0}=F_{1,i}$ and $T_iF_{1,0}=F_{2,ki}$, the latter a positive
integer multiple of $F_{1,i}$. Thus every basis element maps to a positive
integer multiple of a basis element at the next level, not necessarily
to a basis element itself.

Erd\H{o}s and Graham list the irrationality question for $S$~\cite[p.~61]{erdosgraham1980}.

\subsection*{Formal verification and finite certificates}
\paragraph{Formal support.}
The dyadic infinite-basis corollary and the periodic-coefficient corollary
have ordinary proofs above; no additional formal declarations are asserted
for them. For Lemma~\ref{lem:bounded-pulse}, the supplied release records
acceptance by the Lean kernel and by NanoDa at Palomar's Comparator stage.
No new replay was performed for this revision.

\noindent Formal sources: \href{https://github.com/wcook04/plectis-erdos-lean/blob/52f29ad173b04e3bac941b3663f2b9aebe5de0bb/PalomarCorpus/E249/Challenge.lean#L572}{\texttt{wcook04/plectis-erdos-lean}};
\href{https://github.com/wcook04/plectis-erdos-lean/actions/runs/34782407633}{run 34782407633}.

\begin{enumerate}
\item\label{res:rank}
\textbf{Rank in every integer base.} Theorem~\ref{thm:kkernelrank} proves
that the sections through level $e\ge1$ have rank $k^e+1$ for every $k\ge2$.
The formal proof separates the dimension count from the independence
argument. An older theorem assumes independence; the linked
Chinese-remainder and Dirichlet argument proves that assumption.
The resulting rank theorem is unconditional and does not use Martin's
theorem as an axiom. The sources are
\mword{Erdos249257/TotientMahlerDefect.lean}{1084}{finrank_totientKernelThroughLevelFamily_eq}{exact finite-truncation rank};
\mword{Erdos249257/TotientKernelConditional.lean}{215}{finrank_allBaseTotientKernelThroughLevelFamily_eq_of_linearIndependent}{the same count from an independent family};
\href{https://github.com/wcook04/plectis-erdos/blob/25ef6245d15a47548c6926369ae8f1a0f0a14a80/Erdos249257/AllBaseTotientKernel.lean#L1118}{the independence argument in every base};
\href{https://github.com/wcook04/plectis-erdos/blob/25ef6245d15a47548c6926369ae8f1a0f0a14a80/Erdos249257/AllBaseTotientKernel.lean#L1239}{complete truncation rank}.
\end{enumerate}

The long record also gives the exact Farey-interval certificate at
window $K=240$, excluding $S=a/q$ for integers $a$ and
$0<q\le Q_0\approx7.96\times10^{34}$. This finite exclusion is separate
from the quantified irrationality conditions.

\noindent Formal sources: \mword{Erdos249257/TotientMahlerDefect.lean}{1084}{finrank_totientKernelThroughLevelFamily_eq}{dyadic finite-level rank $2^e+1$};
\lword{Erdos249/FullDepthRayAmplifier.lean}{364}{eventually_twoSyndetic_fullDepthKillMultipliers_of_seed}{propagation with depth equal to the shift};
\lword{Erdos249/RankOneSharpFloor.lean}{624}{rankOneSubrankQuotient_sub_theta_two_gt_twentyOne_div_threeTwenty}{lower bound for the rank-one quotients};
\href{https://github.com/wcook04/plectis-erdos/blob/25ef6245d15a47548c6926369ae8f1a0f0a14a80/Erdos249257/AllBaseTotientKernel.lean#L1239}{the unconditional all-base rank theorem};
\mword{Erdos249257/TotientKernelConditional.lean}{215}{finrank_allBaseTotientKernelThroughLevelFamily_eq_of_linearIndependent}{Rank from an independent all-base family};
\lword{Erdos249/CyclotomicAnchoredKill.lean}{3165}{fullMersenneCanonicalBasepointResidueGapSupply_iff_irrational}{the equivalent residue condition}.

\subsection*{Source guide}
Each link retains its repository, commit and line reference. The main
checkpoint is \texttt{\commitshort}; the all-base independence and rank
statements also use \texttt{25ef6245d15a}. Conditional theorems verify
implications, not their unproved hypotheses. The ordinary proofs and
scope statements precede this guide.

\medskip\noindent\textbf{Bases and relations.}\par\smallskip
\mword{Erdos249257/TotientMahlerDefect.lean}{160}{totientKernel_zero_residue}{reduction of sections with residue zero};
\mword{Erdos249257/TotientMahlerDefect.lean}{169}{totientKernel_even_residue_reduce}{reduction of even residues};
\mword{Erdos249257/TotientMahlerDefect.lean}{1265}{linearIndependent_oddCoreTotientKernelFamily}{independence};
\mword{Erdos249257/TotientMahlerDefect.lean}{1380}{span_range_fullTotientKernel_eq_span_range_oddCore}{span equality};
\mword{Erdos249257/TotientMahlerDefect.lean}{1392}{totientDyadicSectionBasis}{basis};
\mword{Erdos249257/TotientMahlerDefect.lean}{935}{linearIndependent_canonicalTotientKernelFamily}{independence};
\mword{Erdos249257/TotientMahlerDefect.lean}{989}{finrank_canonicalTotientKernel_eq}{dimension};
\mword{Erdos249257/TotientMahlerDefect.lean}{1084}{finrank_totientKernelThroughLevelFamily_eq}{the exact finite-truncation rank};
\mword{Erdos249257/TotientMahlerDefect.lean}{1145}{not_finiteDimensional_span_fullTotientKernel}{infinite-dimensionality directly};
\mword{Erdos249257/TotientKernelIndex.lean}{124}{totientKernelSectionIndex_spec}{coordinate specification};
\mword{Erdos249257/TotientKernelIndex.lean}{190}{totientKernelResidueAtLevel_injective}{injective residues};
\mword{Erdos249257/TotientKernelIndex.lean}{213}{existsUnique_totientKernelResidueAtLevel}{unique representation};
\mword{Erdos249257/TotientKernelIndex.lean}{241}{card_totientKernelSectionLevel}{level count};
\mword{Erdos249257/TotientKernelIndex.lean}{249}{card_totientKernelSectionIndex}{section-index count};
\mword{Erdos249257/TotientKernelIndex.lean}{264}{card_totientKernelIndex_eq_two_add_sum}{index as two plus the odd-residue count};
\mword{Erdos249257/TotientKernelIndex.lean}{277}{card_totientKernelIndex}{full index count};
\mword{Erdos249257/TotientMahlerDefect.lean}{823}{paritySeparatedMatrix_det_ne_zero}{a nonzero evaluation determinant};
\mword{Erdos249257/TotientMahlerDefect.lean}{882}{exists_separatedMinorCertificate_totientAffineOddFamily}{an invertible minor for affine totient sections};
\mword{Erdos249257/TotientKernelReduction.lean}{60}{totient_pow_mul_eq}{power-multiplication identity};
\mword{Erdos249257/TotientKernelReduction.lean}{117}{totient_pow_mul_affine_gcd_cross_eq}{affine cross-gcd identity};
\mword{Erdos249257/TotientKernelConditional.lean}{35}{totientKernelReductionScalar_ne_zero}{nonzero reduction scalar};
\mword{Erdos249257/TotientKernelConditional.lean}{54}{allBaseTotientKernelSeq_mul_residue_step}{all-base residue step};
\mword{Erdos249257/TotientKernelConditional.lean}{102}{allBaseTotientKernelSeq_mem_span_canonical_of_le}{canonical spanning};
\mword{Erdos249257/TotientKernelConditional.lean}{185}{span_allBaseTotientKernelThroughLevelFamily_eq_canonical}{span equality};
\mword{Erdos249257/TotientKernelConditional.lean}{203}{finrank_canonicalAllBaseTotientKernel_eq_of_linearIndependent}{rank in every base from independence};
\mword{Erdos249257/TotientKernelConditional.lean}{215}{finrank_allBaseTotientKernelThroughLevelFamily_eq_of_linearIndependent}{all-base truncation rank from independence};
\href{https://github.com/wcook04/plectis-erdos/blob/25ef6245d15a47548c6926369ae8f1a0f0a14a80/Erdos249257/AllBaseTotientKernel.lean#L1118}{independence in every integer base};
\href{https://github.com/wcook04/plectis-erdos/blob/25ef6245d15a47548c6926369ae8f1a0f0a14a80/Erdos249257/AllBaseTotientKernel.lean#L1231}{the rank of the retained all-base family};
\href{https://github.com/wcook04/plectis-erdos/blob/25ef6245d15a47548c6926369ae8f1a0f0a14a80/Erdos249257/AllBaseTotientKernel.lean#L1239}{the rank of all sections through a finite level}.

\medskip\noindent\textbf{Tail identities, residue tests and carries.}\par\smallskip
\mword{Erdos249257/TotientTailPeriodKiller.lean}{150}{two_pow_mul_totient_series_eq}{the finite-prefix and tail identity};
\href{https://github.com/wcook04/plectis-erdos/blob/99f4bf47422abbd8757cbb22b50ba079d764d3a7/Erdos249257/TotientTailPeriodKiller.lean#L72}{definition};
\mword{Erdos249257/TotientTailPeriodKiller.lean}{404}{certifiedKill_all_small}{certificates for shifts 1 through 8};
\mword{Erdos249257/TotientTailPeriodKiller.lean}{407}{tail_diff_not_int_small}{nonintegral tail differences at the checked shifts};
\mword{Erdos249257/GenericTailOrbitRigidity.lean}{67}{IsTemperedBinaryOrbit}{the integer carry recurrence and its growth bound};
\lword{Erdos249/CyclotomicAnchoredKill.lean}{2322}{fullMersenneBlockResidue_succ}{residue recurrence};
\lword{Erdos249/CyclotomicAnchoredKill.lean}{3132}{fullMersenneCenteredResidueGapSupply_of_canonicalBasepoint}{comparison of the two residue conditions};
\lword{Erdos249/CyclotomicAnchoredKill.lean}{3165}{fullMersenneCanonicalBasepointResidueGapSupply_iff_irrational}{exact equivalence};
\lword{Erdos249/FullDepthRayAmplifier.lean}{364}{eventually_twoSyndetic_fullDepthKillMultipliers_of_seed}{a certificate in each sufficiently late pair of multipliers};
\lword{Erdos249/FullDepthRayAmplifier.lean}{388}{exists_fullDepthKill_on_ray_iff_shift_notMem_int}{pointwise equivalence};
\lword{Erdos249/FullDepthRayAmplifier.lean}{406}{apFullDepthEscape_iff_irrational}{global equivalence};
\lword{Erdos249/FullDepthRayAmplifier.lean}{421}{cofinalFullDepthKillSupply_iff_periodMultipleKillSupply}{equivalence after setting depth equal to the shift};
\lword{Erdos249/FullDepthRayAmplifier.lean}{438}{cofinalFullDepthKillSupply_iff_irrational}{irrationality and arbitrarily large certificates};
\mword{Erdos249257/TotientCarryKernelRigidity.lean}{300}{not_irrational_totientSeries_implies_unbounded_carryRank_unconditional}{the carry-rank lower bound under rationality};
\mword{Erdos249257/TotientCarryKernelRigidity.lean}{211}{finrank_canonicalCarryKernel_ge_of_linearIndependent}{rank transport from \ref{res:rank}};
\mword{Erdos249257/TotientTailCarryPeriod.lean}{224}{not_irrational_totientSeries_implies_mod_period_and_unbounded_rank}{modular periodicity and the rank lower bound for the same carry};
\mword{Erdos249257/TotientTailCarryPeriod.lean}{126}{totient_forcing_vanishes_mod_of_dvd_multiplier}{forcing vanishing};
\mword{Erdos249257/TotientTailCarryPeriod.lean}{140}{totient_carry_modEq_geometric_of_dvd_multiplier}{geometric reduction};
\mword{Erdos249257/CertificateKernel.lean}{18445}{tsum_primWeight_div_two_pow_sub_one_eq_totient_series}{identity};
\mword{Erdos249257/MersenneLambertLadder.lean}{370}{primWeight_nonneg}{sign};
\mword{Erdos249257/MersenneLambertLadder.lean}{297}{primWeight_apply_prime}{primes};
\mword{Erdos249257/MersenneLambertLadder.lean}{272}{primWeight_apply_prime_pow}{prime powers};
\mword{Erdos249257/MersenneLambertLadder.lean}{321}{primWeight_not_bounded}{unboundedness};
\mword{Erdos249257/SquaredMersenneDiagonalEnclosure.lean}{138}{scaleDiagonalTailDifference_sub_lambertProjectedDiagonal}{exact centre decomposition};
\mword{Erdos249257/SquaredMersenneDiagonalEnclosure.lean}{408}{abs_scaleDiagonalTailDifference_sub_lambertProjectedDiagonal_le}{explicit error bound};
\mword{Erdos249257/SquaredMersenneDiagonalEnclosure.lean}{426}{scaleFullTarget_miss_of_lambert_projected_separation}{separation criterion};
\mword{Erdos249257/FullTargetPrimeAdjunctionNoGo.lean}{139}{scaleFullTargetHit_iff_integral}{target equivalence};
\mword{Erdos249257/FreshPrimeDeficitDecomposition.lean}{91}{endpointFreshDeficit_nonneg}{endpoint nonnegativity};
\mword{Erdos249257/FreshPrimeDeficitDecomposition.lean}{184}{finiteForeignChannelIncrement_eq_endpointFreshDeficit_sub}{foreign-channel identity};
\mword{Erdos249257/FreshPrimeDeficitDecomposition.lean}{269}{oldFivePoint_sub_secondBranchAdverseDeficit_le_actual}{five-point loss bound};
\mword{Erdos249257/SquareCRTCube.lean}{297}{exists_squareCRT_clean_totient_family}{simultaneous finite family};
\mword{Erdos249257/SquareCRTCube.lean}{327}{exists_squareCRT_clean_horizon}{finite horizon};
\mword{Erdos249257/SquareCRTCube.lean}{459}{squareCRT_clean_block_can_vanish}{vanishing block};
\mword{Erdos249257/SquareCRTCube.lean}{477}{squareCRT_clean_block_can_be_nonzero}{nonzero block};
\mword{Erdos249257/DiagonalFreshLossBridge.lean}{2667}{canonicalAdjacentSuffixCentral_iff_slack_nonneg}{centrality equivalence};
\mword{Erdos249257/DiagonalFreshLossBridge.lean}{2757}{canonicalAdjacentSuffixJumpCentralSupply_iff_slackSupply}{equivalence of the two conditions};
\mword{Erdos249257/DiagonalFreshLossBridge.lean}{2932}{irrational_totientSeries_of_canonicalAdjacentSuffixJumpSlackSupply}{conditional irrationality theorem};
\mword{Erdos249257/ActualForeignResidueProjection.lean}{276}{abs_foreignTailWindow_le_foreignComplementBound}{finite complement bound};
\mword{Erdos249257/ActualForeignResidueProjection.lean}{308}{finiteResidueDiagonal_eq_projectedForeign_add_divisor}{exact channel split};
\mword{Erdos249257/ActualForeignResidueProjection.lean}{414}{scaleFullTarget_miss_of_projected_separation}{conditional target-miss theorem};
\mword{Erdos249257/CertificateKernel.lean}{12811}{irrational_ratWeightSeries_eventuallyPeriodic}{the periodic theorem};
\mword{Erdos249257/CertificateKernel.lean}{18429}{tsum_totient_div_two_pow_sub_one_eq_two}{$\sum\ph(n)/(2^n-1)=2$};
\mword{Erdos249257/GenericTailOrbitRigidity.lean}{426}{binaryCoeffSeries_rational_iff_exists_temperedBinaryOrbit}{rationality and integer carries with controlled growth};
\mword{Erdos249257/FreshPrimeDeficitDecomposition.lean}{171}{diagonalHeightIncrement_eq_old_add_endpointFreshDeficits}{diagonal split};
\href{https://github.com/wcook04/plectis-erdos/blob/99f4bf47422abbd8757cbb22b50ba079d764d3a7/Erdos249257/CarrySurvivorExtinction.lean#L515}{definition};
\mword{Erdos249257/TotientTailPeriodKiller.lean}{57}{totientTail}{definition of the scaled tail};
\lword{Erdos249/CyclotomicAnchoredKill.lean}{3231}{CyclotomicAnchoredKillSupply}{the proposed cyclotomic residue condition};
\lword{Erdos249/CyclotomicAnchoredKill.lean}{3327}{binaryCyclotomicAnchoredKillSupply_iff_irrational}{the checked equivalence};
\lword{Erdos249/CyclotomicAnchoredKill.lean}{1706}{exists_unbounded_binaryCyclotomicSupport_with_periodLock_of_not_irrational}{checked conditional support/period counterexample};
\lword{Erdos249/CyclotomicAnchoredKill.lean}{1693}{binaryCyclotomicLayer_unboundedPrimeDivisorSupply}{the checked binary-layer theorem};
\lword{Erdos249/PeriodMultipleEscape.lean}{441}{ApFullDepthEscape}{certificates with depth equal to the period multiple}.

\medskip\noindent\textbf{Finite denominator exclusions.}\par\smallskip
\mword{Erdos249257/CertificateKernel.lean}{18371}{tsum_totient_div_pow_two_ne_int_div_of_den_le_79639646646701375323355774875831053}{exclusion};
\mword{Erdos249257/CertificateKernel.lean}{18384}{tsum_totient_div_pow_two_ne_ratCast_of_den_le_79639646646701375323355774875831053}{reduced-denominator form};
\mword{Erdos249257/GapFareyBound.lean}{225}{gap_check_window_1_240_first_failure}{first failure};
\mword{Erdos249257/CertificateKernel.lean}{16025}{irrational_tsum_totient_div_pow_two_of_gap_certificate_supply}{the conditional implication};
\mword{Erdos249257/GapFareyBound.lean}{51}{farey_gap}{farey gap};
\mword{Erdos249257/GapFareyBound.lean}{176}{gap_check_window_1_240_le_79639646646701375323355774875831053}{the denominator bound at window 240};
\mword{Erdos249257/CertificateKernel.lean}{15986}{tsum_totient_div_pow_two_ne_int_div_of_totient_gap_certificate}{denominator exclusion from a finite interval};
\mword{Erdos249257/GapFareyBound.lean}{88}{gap_check_window_1_120_le_248672326362367909}{the earlier denominator bound at window 120};
\mword{Erdos249257/CertificateKernel.lean}{18572}{tsum_visible_coprime_pairs_ne_int_div_of_den_le_39819823323350687661677887437915526}{denominator exclusion for the coprimality probability};
\mword{Erdos249257/DiagonalPincerCertificatesT64.lean}{1933}{diagonalPincerCertificateScalesThroughT64}{list};
\mword{Erdos249257/DiagonalPincerCertificatesT64.lean}{1967}{certifiedKill_diagonal_all_imported_through_t64}{verification};
\lword{Skip/LadderT67.lean}{71264}{exists_diagonalKill_le_82}{certificates at every scale through 82}.

\medskip\noindent\textbf{Least common multiples and limits of the tested methods.}\par\smallskip
\mword{Erdos249257/LcmConeFlatness.lean}{316}{exists_certifiedKill_iff_tail_diff_notMem_int}{certificate completeness};
\mword{Erdos249257/LcmConeFlatness.lean}{327}{tail_diff_mem_int_iff_scaled_series_mem_int}{tail-difference integrality};
\mword{Erdos249257/LcmConeFlatness.lean}{386}{irrational_totient_series_iff_all_tail_diffs_nonintegral}{irrationality iff every tail difference is nonintegral};
\mword{Erdos249257/TotientActualLcmOrbitNonintegrality.lean}{37}{irrational_totientSeries_iff_actualLcmOrbitNonintegralitySupply}{exact frontier};
\mword{Erdos249257/TotientActualLcmOrbitNonintegrality.lean}{53}{irrational_totientSeries_of_actualLcmOrbitNonintegralitySupply}{forward implication};
\mword{Erdos249257/TotientActualLcmOrbitSign.lean}{39}{actualLcmTailDiff_shift_pos}{positive tail differences};
\mword{Erdos249257/TotientActualLcmOrbitSign.lean}{172}{actualLcm_trueEndpointSurvivor_neg}{negative endpoint representative};
\mword{Erdos249257/TotientActualLcmOrbitSign.lean}{211}{actualLcm_integral_forces_topEdgeResidue}{residue near the upper endpoint};
\mword{Erdos249257/LcmFactorIdealPulseObstruction.lean}{798}{lcm_factorIdeal_finiteRank_shiftAlgebra_not_sufficient}{divisibility preserved by finite combinations of shifts};
\mword{Erdos249257/LcmFactorIdealPulseObstruction.lean}{866}{lcm_factorIdeal_sparseAnchor_not_sufficient}{the counterexample with sparse prescribed values};
\mword{Erdos249257/LcmConeFlatness.lean}{399}{irrational_totient_series_iff_pointwise_certificates}{pointwise};
\mword{Erdos249257/LcmConeFlatness.lean}{412}{irrational_totient_series_iff_certificate_supply}{cofinal at each shift};
\mword{Erdos249257/LcmConeFlatness.lean}{426}{irrational_totient_series_iff_lcm_diagonal_certificate_supply}{$\lcm$ diagonal};
\mword{Erdos249257/TotientActualLcmShortKill.lean}{18}{lcmDiagonalArithmeticKill_two_pow_four}{the certificate at \(t=2^4\) and depth \(23\)};
\mword{Erdos249257/TotientActualLcmShortKill.lean}{27}{lcmDiagonalArithmeticKill_two_pow_six}{the \(t=2^6\), length-\(93\) kill};
\mword{Erdos249257/TotientActualLcmShortKill.lean}{35}{actualLcmTailOrbit_four_and_six_notMem_int}{actual-orbit non-integrality at both heights};
\mword{Erdos249257/TotientActualLcmShortKill.lean}{54}{powerTwoActualLcmShortArithmeticKillSupply_through_six}{but explicitly stops at that finite prefix};
\mword{Erdos249257/TotientActualLcmOrbitSeparation.lean}{141}{abs_actualLcmTailOrbit_sub_rawApprox_lt}{recorded here};
\mword{Erdos249257/TotientActualLcmOrbitSeparation.lean}{179}{actualLcmRawApprox_eq_half_div_fourPow}{by the normalization identity};
\mword{Erdos249257/TotientActualLcmOrbitSeparation.lean}{208}{halfWordBandAt_of_rawApprox_integerSeparation}{through the exact threshold lemma};
\mword{Erdos249257/TotientActualLcmOrbitSeparation.lean}{254}{PowerTwoActualLcmOrbitSeparationSupply}{the cofinal separation hypothesis};
\mword{Erdos249257/TotientActualLcmOrbitSeparation.lean}{305}{irrational_totientSeries_of_actualLcmOrbitSeparationSupply}{the conditional endpoint theorem};
\mword{Erdos249257/TotientActualLcmTopEdgeStaircase.lean}{251}{not_actualLcmTerminalDyadicStaircase_of_room}{the staircase is impossible under the stated room hypothesis};
\mword{Erdos249257/TotientActualLcmTopEdgeStaircase.lean}{297}{windowDiscrepancy_emod_two_pow_eq_terminal}{terminal-remainder identity};
\mword{Erdos249257/TropicalCurvatureCarry.lean}{137}{fixedPrecisionTropicalNoGo}{the limitation of fixed-precision valuation data};
\mword{Erdos249257/LcmConeFlatness.lean}{357}{not_irrational_of_tail_diff_mem_int}{rationality from an integral tail difference}.

\medskip\noindent\textbf{M\"obius identities and rational approximation.}\par\smallskip
\lword{Erdos249/RankOneSharpFloor.lean}{534}{rankOneSubrankQuotient_eq_one_five_iff}{unique minimizer};
\lword{Erdos249/RankOneSharpFloor.lean}{624}{rankOneSubrankQuotient_sub_theta_two_gt_twentyOne_div_threeTwenty}{uniform lower bound $21/320$};
\lword{Erdos249/RankOneSharpFloor.lean}{634}{rankOneSubrankQuotient_sub_theta_two_gt_one_div_sixteen}{lower bound $1/16$};
\lword{Erdos249/RankOneSharpFloor.lean}{645}{not_forall_rankOneSubrankQuotient_sub_theta_two_gt_one_div_fifteen}{$1/15$ failure};
\lword{Erdos249/RankOneSharpFloor.lean}{655}{positive_direct_sum_sub_theta_two_gt_twentyOne_div_threeTwenty}{positive weighted averages of the quotients};
\lword{Erdos249/RankOneSharpFloor.lean}{693}{primitive_form_abs_gt_twentyOne_div_threeTwenty}{rational linear-form obstruction};
\mword{Erdos249257/GcdMomentCalculus.lean}{216}{tsum_one_div_mersenne_sq_eq_sigma_sub_tau_series}{divisor-count layer};
\mword{Erdos249257/GcdMomentCalculus.lean}{235}{tsum_totient_div_mersenne_sq_eq_gcd_moment_series}{first gcd moment};
\mword{Erdos249257/PrimitiveDeterminantLift.lean}{169}{no_fixed_integral_primitive_euler_index}{no fixed index};
\mword{Erdos249257/PrimitiveDeterminantLift.lean}{148}{twoTierPrimorial_dvd_of_integral_jet}{the primorial divisibility};
\mword{Erdos249257/CertificateKernel.lean}{8335}{irrational_erdosBorwein_series}{irrationality of the Erd\H{o}s--Borwein series};
\mword{Erdos249257/CertificateKernel.lean}{18434}{tsum_moebius_div_two_pow_sub_one_eq_half}{$\sum\mu(n)/(2^n-1)=1/2$};
\mword{Erdos249257/CertificateKernel.lean}{18454}{totient_series_eq_half_add_moebius_mersenne_square}{the squared-Mersenne expression for the totient series};
\mword{Erdos249257/CertificateKernel.lean}{18557}{totient_series_eq_half_add_visible_coprime_pairs}{the coprimality-probability expression};
\mword{Erdos249257/GcdMomentCalculus.lean}{349}{tsum_pos_coprime_inv_mersenne_eq_one}{the Lambert-weighted sum over coprime positive pairs};
\mword{Erdos249257/GcdMomentCalculus.lean}{474}{cylinderMass_split}{the Stern--Brocot recursion with stopping};
\lword{Erdos249/FiniteEulerSieve.lean}{28}{muSqEulerFactor_one}{$s=1$};
\lword{Erdos249/FiniteEulerSieve.lean}{36}{muSqEulerFactor_two}{$s=2$};
\lword{Erdos249/FiniteEulerSieve.lean}{44}{sigmaPrimePow_firstDiff}{first difference};
\lword{Erdos249/FiniteEulerSieve.lean}{51}{sigmaPrimePow_secondDiff}{second difference}.

\medskip\noindent\textbf{Proposed exponential-sum and prime-factor estimates.}\par\smallskip
\mword{Erdos249257/PivotAntiReconstruction.lean}{1765}{dtwWindowSeparatedPairs_iff_irrational_totient_series}{window-separated pairs};
\lword{Erdos249/TotientStrictPrimeEscape.lean}{157}{DTWNaturalPrimeTailOrbitStrictGap}{the unproved quantitative separation condition at primes};
\lword{Erdos249/TotientStrictPrimeEscape.lean}{485}{naturalPivotPointEscape_of_naturalPrimeTailOrbitStrictGap}{a positive adaptive truncation budget};
\lword{Erdos249/TotientStrictPrimeEscape.lean}{524}{irrational_totient_series_of_naturalPrimeTailOrbitStrictGap}{$S$ is irrational};
\href{https://github.com/wcook04/plectis-erdos/blob/99f4bf47422abbd8757cbb22b50ba079d764d3a7/ErdosProblems/Erdos249/TotientStrictPrimeEscape.lean#L35}{density hypothesis};
\href{https://github.com/wcook04/plectis-erdos/blob/99f4bf47422abbd8757cbb22b50ba079d764d3a7/ErdosProblems/Erdos249/TotientStrictPrimeEscape.lean#L25}{block-gap hypothesis};
\href{https://github.com/wcook04/plectis-erdos/blob/99f4bf47422abbd8757cbb22b50ba079d764d3a7/ErdosProblems/Erdos249/TotientStrictPrimeEscape.lean#L45}{density-to-gap implication};
\href{https://github.com/wcook04/plectis-erdos/blob/99f4bf47422abbd8757cbb22b50ba079d764d3a7/ErdosProblems/Erdos249/TotientStrictPrimeEscape.lean#L167}{exact squaring recurrence};
\href{https://github.com/wcook04/plectis-erdos/blob/99f4bf47422abbd8757cbb22b50ba079d764d3a7/ErdosProblems/Erdos249/TotientStrictPrimeEscape.lean#L191}{iterate formula};
\href{https://github.com/wcook04/plectis-erdos/blob/99f4bf47422abbd8757cbb22b50ba079d764d3a7/ErdosProblems/Erdos249/TotientStrictPrimeEscape.lean#L201}{initial-phase normal form};
\href{https://github.com/wcook04/plectis-erdos/blob/99f4bf47422abbd8757cbb22b50ba079d764d3a7/ErdosProblems/Erdos249/TotientStrictPrimeEscape.lean#L278}{initial-phase equivalence};
\href{https://github.com/wcook04/plectis-erdos/blob/99f4bf47422abbd8757cbb22b50ba079d764d3a7/ErdosProblems/Erdos249/TotientStrictPrimeEscape.lean#L322}{dyadic-root obstruction};
\href{https://github.com/wcook04/plectis-erdos/blob/99f4bf47422abbd8757cbb22b50ba079d764d3a7/ErdosProblems/Erdos249/TotientStrictPrimeEscape.lean#L465}{prime-alignment implication};
\lword{Erdos249/PrimeRayCyclotomicCurvature.lean}{22}{checkerboard}{mixed difference};
\lword{Erdos249/PrimeRayCyclotomicCurvature.lean}{26}{checkerboard_separable}{separable case};
\lword{Erdos249/PrimeRayCyclotomicCurvature.lean}{32}{checkerboard_unique}{fixed-stencil uniqueness};
\lword{Erdos249/PrimeRayCyclotomicCurvature.lean}{45}{fourPoint_layer_identity}{a four-point divisor-layer identity};
\lword{Erdos249/PrimeRayCyclotomicCurvature.lean}{62}{residueClass_step_has_centred_completion}{a centred representative for a residue class modulo an even modulus};
\lword{Erdos249/PrimeRayCyclotomicCurvature.lean}{95}{PrimeRayLayerSupply}{eventual nontriviality and coprimality of the layers};
\lword{Erdos249/PrimeRayCyclotomicCurvature.lean}{102}{BoundedDegreeOrderConsumer}{order divisibility with bounded extension degree};
\lword{Erdos249/PrimeRayCyclotomicCurvature.lean}{151}{FinitePrimeSupportEscape}{exclusion of any fixed finite prime support};
\lword{Erdos249/PrimeRayCyclotomicCurvature.lean}{158}{UnboundedPrimeDivisorSupply}{unbounded prime divisors};
\mword{Erdos249257/FirstHarmonicPivot.lean}{299}{mem_pivotFiber_iff_exists_supplierPrime}{membership};
\mword{Erdos249257/FirstHarmonicPivot.lean}{367}{pivotBaseOfPrime_injective_on_supplierPrimes}{injectivity};
\mword{Erdos249257/FirstHarmonicPivot.lean}{384}{image_pivotSupplierPrimes_eq_pivotFiber}{image equality};
\mword{Erdos249257/FirstHarmonicPivot.lean}{394}{supplierPrime_not_globally_isolated_counterexample}{checked counterexample};
\mword{Erdos249257/FirstHarmonicPivot.lean}{514}{windowFirstExp_sum_eq_pivot_decomposition}{the exact decomposition into four prime-factor contributions};
\mword{Erdos249257/FirstHarmonicPivot.lean}{577}{irrational_totient_series_of_pivotResidualDecorrelation}{the conditional irrationality theorem};
\mword{Erdos249257/FirstHarmonicPivot.lean}{543}{PivotBudgetAt}{the four inequalities on a common block};
\mword{Erdos249257/FirstHarmonicPivot.lean}{549}{first_harmonic_gap_of_pivotBudgetAt}{the first-harmonic bound from those four inequalities}.

\begin{thebibliography}{99}\small
\bibitem{allouche-shallit} J.-P.~Allouche and J.~Shallit,
\emph{The ring of $k$-regular sequences}, Theoret.\ Comput.\ Sci.\
\textbf{98} (1992), no.~2, 163--197,
doi:\href{https://doi.org/10.1016/0304-3975(92)90001-V}{10.1016/0304-3975(92)90001-V}.
\bibitem{coons} M.~Coons, \emph{(Non)Automaticity of number theoretic
functions}, J.\ Th\'eor.\ Nombres Bordeaux \textbf{22} (2010), no.~2, 339--352,
doi:\href{https://doi.org/10.5802/jtnb.718}{10.5802/jtnb.718}.
Theorem~3.2, p.~348.
\bibitem{coons-jsr} M.~Coons, \emph{Regular sequences and the joint spectral
radius}, Internat.\ J.\ Found.\ Comput.\ Sci.\ \textbf{28} (2017), no.~2,
135--140,
doi:\href{https://doi.org/10.1142/S0129054117500095}{10.1142/S0129054117500095};
arXiv:\href{https://arxiv.org/abs/1511.07535v1}{1511.07535v1}.
Theorem~1, Proposition~4, Corollary~7 and Appendix~A are cited with the
arXiv v1 numbering.
\bibitem{martin-phi-inequalities} G.~Martin, \emph{Simultaneous inequalities
among values of the Euler phi-function},
arXiv:\href{https://arxiv.org/abs/math/0603053v1}{math/0603053v1}, 2006.
Theorem~1, pp.~1--2.
\bibitem{wong2015} Erick Wong,
\href{https://math.stackexchange.com/a/1211557}{answer 1211557}
to \emph{An infinite sum based on the mod-parity of Euler's totient function},
Mathematics Stack Exchange, 29 March 2015, accessed 15 September 2026.
\bibitem{eldar2026oeis} A.~Eldar, comment and formula added to OEIS A256936
(revisions 28 and 31), 15 March 2026,
\url{https://oeis.org/history?seq=A256936}, accessed 16 September 2026.
\bibitem{fan2026totient} S.~Fan, comment on Erd\H{o}s Problem \#249,
16 May 2026, 19:01, \href{https://www.erdosproblems.com/forum/thread/249}
{Erd\H{o}s Problems discussion thread}, accessed 16 September 2026.
\bibitem{erdosgraham1980} P.~Erd\H{o}s and R.~L.~Graham,
\emph{Old and New Problems and Results in Combinatorial Number Theory},
Monographies de L'Enseignement Math\'ematique \textbf{28}, 1980, p.~61.
\bibitem{bgs2016} R.~Balasubramanian, S.~Giri and P.~Srivastav,
\emph{On correlations of certain multiplicative functions},
J.\ Number Theory \textbf{174} (2017), 221--238,
doi:\href{https://doi.org/10.1016/j.jnt.2016.10.001}{10.1016/j.jnt.2016.10.001};
arXiv:\href{https://arxiv.org/abs/1511.02221v3}{1511.02221v3}.
Theorem~2.2 is cited with the arXiv v3 pagination, p.~2.
\bibitem{yazdani2001} S.~Yazdani,
\emph{Multiplicative functions and $k$-automatic sequences},
J.\ Th\'eor.\ Nombres Bordeaux \textbf{13} (2001), no.~2, 651--658,
\url{https://www.numdam.org/item/JTNB_2001__13_2_651_0/}.
Theorem~2 and its proof, pp.~652--653; Corollary~4, p.~654.
\bibitem{allouche-shallit-yassawi} J.-P.~Allouche, J.~Shallit and R.~Yassawi,
\emph{How to prove that a sequence is not automatic},
Exposition.\ Math.\ \textbf{40} (2022), no.~1, 1--22,
doi:\href{https://doi.org/10.1016/j.exmath.2021.08.001}{10.1016/j.exmath.2021.08.001};
arXiv:\href{https://arxiv.org/abs/2104.13072v1}{2104.13072v1}.
Theorem and example numbering follows the arXiv v1 (2021).
\bibitem{bell-smertnig2026} J.~Bell and D.~Smertnig,
\emph{Mahler series with multiplicative coefficient sequences},
arXiv:\href{https://arxiv.org/abs/2603.23456v1}{2603.23456v1},
24 March 2026, preprint. Theorem~1.3, p.~2; Lemma~3.3, p.~9.
\bibitem{bell-bruin-coons} J.~P.~Bell, N.~Bruin and M.~Coons,
\emph{Transcendence of generating functions whose coefficients are multiplicative},
Trans.\ Amer.\ Math.\ Soc.\ \textbf{364} (2012), no.~2, 933--959,
doi:\href{https://doi.org/10.1090/S0002-9947-2011-05479-6}{10.1090/S0002-9947-2011-05479-6};
arXiv:\href{https://arxiv.org/abs/1003.2221v2}{1003.2221v2}.
Theorems~1.5--1.6 are cited with the arXiv v2 pagination, p.~2.
\bibitem{adamczewski-bell-smertnig2023}
B.~Adamczewski, J.~Bell and D.~Smertnig,
\emph{A height gap theorem for coefficients of Mahler functions},
J.~Eur.~Math.~Soc.\ \textbf{25} (2023), no.~7, 2525--2571,
doi:\href{https://doi.org/10.4171/JEMS/1244}{10.4171/JEMS/1244}.
Theorem~1.2(b), p.~2528; logarithmic height in Section~1.1.1, p.~2527.
\end{thebibliography}

\companionpapercontext

\end{document}
